% cap3.tex

\chapter{Desarrollo Metodológico} \label{MCP:subastador} % la etiqueta para referencias

Este capítulo describe el plan de trabajo de \textit{Proyecto de Título II}
(PT2, segundo semestre de 2026), acordado con el profesor guía Sebastián Cea.
Cruza la guía de trabajo del semestre (\texttt{ESTRATEGIA\_Semestre2.md}) con la
propuesta de hitos entregada el 21 de agosto de 2026
(\texttt{Propuesta\_de\_Hitos\_PT2.tex}), e incorpora la actualización acordada
con el profesor guía en una reunión posterior a ambos documentos, que redefine
el orden de los dos hitos de avance del semestre.

\section{Objetivo general de PT2}

El resultado central del modelo de \citeA{venegas_reyes_financial_2025}
—recuperación completa del bienestar mediante una infraestructura alternativa,
frente a un \textit{gap} de $27{,}9\%$ que deja la sola competencia entre
intermediarios— se obtuvo bajo agentes homogéneos y parámetros ilustrativos.
El objetivo de PT2 es producir el insumo de \textbf{validación empírica} que la
tesis de Venegas (2024) no desarrolla: evaluar si ese resultado es consistente
con datos reales, calibrando los primitivos del modelo ($\mu_i$, $\Sigma_i$,
$\gamma_i$) contra un caso concreto de datos públicos de financiamiento en
Chile, en lugar de valores ilustrativos.

A diferencia de lo planteado en \texttt{Hito\_2.tex} durante PT1, este semestre
\textbf{no se aborda la extensión del modelo a agentes heterogéneos} en
aversión al riesgo y creencias. Esa extensión queda como línea de trabajo
futura, pendiente de una decisión posterior con el profesor guía, y las tres
preguntas de investigación heredadas de Hito 2 se trabajan en su versión
homogénea (Capítulo~\ref{c1}).

\section{Actualización del orden de hitos}

\texttt{ESTRATEGIA\_Semestre2.md} y \texttt{Propuesta\_de\_Hitos\_PT2.tex}
—ambos redactados antes de esta actualización— describen el hito del 27 de
septiembre como la formalización analítica del modelo heterogéneo y su
validación numérica contra los \textit{benchmarks} de Venegas, dejando la
aplicación a datos reales para el hito del 8 de noviembre. En una reunión
posterior con el profesor guía, este orden se invirtió: el primer hito de
avance pasa a ser una \textbf{prueba de factibilidad empírica}, y el segundo,
la calibración y aplicación completas. Este capítulo refleja esa actualización;
los dos documentos de planificación citados arriba aún no se han corregido en
el repositorio para reflejarla, lo que el estudiante debe confirmar
explícitamente con el profesor guía antes de la entrega de septiembre.

\section{Hito I — 27 de septiembre de 2026: prueba de factibilidad empírica}

Evaluar qué tan factible es tomar el modelo de \citeA{venegas_reyes_financial_2025}
—en su versión homogénea, sin extensión a heterogeneidad— y cruzarlo con datos
reales. Comprende:

\begin{itemize}
    \item Identificación y evaluación de acceso a fuentes de datos públicos de
    financiamiento en Chile, en particular a través del Banco Interamericano de
    Desarrollo (BID) y de CORFO, que permitan aproximar los primitivos del
    modelo ($\mu_i$, $\Sigma_i$, $\gamma_i$) para un conjunto de agentes reales
    (p. ej. bancos y empresas beneficiarias de líneas de financiamiento CORFO).
    \item Verificación de consistencia numérica del modelo base implementado
    (\texttt{code/Financial\_Networks\_Stability.qmd}) contra los
    \textit{benchmarks} ya publicados por Venegas (2024): escenarios
    3F-WITH/WITHOUT y 4F-WITH/WITHOUT.
    \item Un primer diagnóstico de factibilidad: si los datos disponibles
    permiten una calibración razonable de los primitivos del modelo, o si se
    requieren supuestos adicionales o un caso alternativo.
\end{itemize}

\subsection{Diagnóstico muestral preliminar}

La Sección~\ref{c25:identificacion} estableció que estimar $\Sigma$ mediante la
Proposición~1 de \citeA{jalan_incentive-aware_2024} requiere una secuencia
longitudinal $W(t)$ observada en $T$ periodos, y que para una red de $n=3$
agentes el Teorema~5 sugiere un mínimo del orden de $4$ a $5$ observaciones por
agente. Un primer levantamiento de datos públicos sobre el programa
\textit{Crédito Corfo Mipyme}, a partir de las Cuentas Públicas Participativas
de CORFO, permitió construir un panel preliminar de flujos anuales
BID/CORFO$\to$intermediaria ($W_{12}$) e intermediaria$\to$MiPyME ($W_{23}$)
para los años 2016, 2019, 2020, 2021, 2023 y 2024 ($T=6$ periodos con cifras
utilizables, con 2014, 2015, 2017, 2018 y 2022 aún sin cifra nacional
comparable localizada). Aun con esos vacíos, $T=6$ ya satisface el mínimo
teórico de la Proposición~1 para $n=3$, lo que sostiene un diagnóstico de
factibilidad muestral positivo para el Hito~I.

Dos advertencias matizan esa conclusión preliminar. Primero, $W_{23}$ no es un
flujo derivado únicamente del $W_{12}$ del mismo año —incluye capital rotativo
de desembolsos previos que siguen vigentes—, lo que genera una tensión
conceptual con el modelo estático de \citeA{jalan_incentive-aware_2024} y
\citeA{venegas_reyes_financial_2025}, que tratan $W_{ij}$ como el tamaño de un
contrato vigente en un punto estable, no como un flujo acumulado; queda
pendiente decidir si $W_{12}(t)$ debe reconstruirse como saldo vigente o si se
acepta $W_{23}$ como aproximación de flujo anual. Segundo, la arista prohibida
BID/CORFO--MiPyME de esta red implica que la entrada $\Sigma_{13}$ queda débilmente
identificada por la Proposición~1 (Sección~\ref{c25:identificacion}), por lo
que la calibración completa del Hito~II probablemente deba reportar $\Sigma_{13}$
como no identificado en vez de estimarlo directamente.

\section{Hito II — 8 de noviembre de 2026: calibración y aplicación}

Con el diagnóstico de factibilidad del Hito I resuelto, este hito comprende:

\begin{itemize}
    \item Calibración completa del modelo con los datos reales seleccionados.
    \item Aplicación del modelo calibrado para contrastar sus predicciones de
    recuperación de bienestar y de reparto de actividad entre la red
    tradicional y la alternativa contra la evidencia observada en el caso real.
    \item Análisis de sensibilidad de los resultados ante variaciones en los
    parámetros calibrados ($\tau$, $\lambda_A$).
\end{itemize}

\section{Hitos de cierre — 21 de agosto y 28 de diciembre de 2026}

El 21 de agosto de 2026 se entregó la \texttt{Propuesta\_de\_Hitos\_PT2.tex},
con la carta Gantt del semestre re-secuenciada contra las cuatro fechas de este
plan. El 28 de diciembre de 2026 corresponde a la entrega del documento de
memoria completo (\texttt{Caballero.tex}): los cinco capítulos redactados y
compilando sin errores, los resultados de la aplicación a datos reales
integrados, y \texttt{referencias.bib} consolidado.
